% 原章第44—54页：5.5.3、软件包；式5.110—5.138及结论5.10。
\subsection{内生变量与外生变量均为函数的函数型结构方程模型（FSEMF）}
\subsubsection{模型}

现在考虑 meth-seq 与下一代测序（NGS）基因型数据。基因型—甲基化网络既连接甲基化基因与甲基化基因，也将含有密集分布 SNP 的基因连接到甲基化基因。假定有 $M$ 个甲基化基因、$K$ 个经过测序且包含罕见变异或同时包含常见与罕见变异的基因，并抽取了 $T$ 名个体。令 $y_m(s_m)$（$m=1,\ldots,M$）为表示基因组位置 $s_m$ 处 CpG 位点甲基化水平的变量，$y_{im}(s_m)$ 为第 $i$ 名个体所取的值。将 $y_m(s_m)$ 定义为甲基化函数。令 $x_k(t)$ 为位于基因组位置 $t$ 处 SNP 基因型的指示变量。假定全部 $M$ 个甲基化函数与 $K$ 个基因型函数都映射到 $[0,1]$ 区间。

定义 FSEMF：
\begin{equation}
\begin{aligned}
&y_1(s)r_{11}+\cdots+y_M(s)r_{M1}\\
&\quad+\int_Tx_1(t)\beta_{11}(s,t)\,dt+\cdots
+\int_Tx_K(t)\beta_{K1}(s,t)\,dt+e_1(s)=0,\\
&\hspace{14em}\vdots\\
&y_1(s)r_{1M}+\cdots+y_M(s)r_{MM}\\
&\quad+\int_Tx_1(t)\beta_{1M}(s,t)\,dt+\cdots
+\int_Tx_K(t)\beta_{KM}(s,t)\,dt+e_M(s)=0,
\end{aligned}
\tag{5.110}
\end{equation}
其中，$s\in[0,1]$ 为映射后的基因组位置；$\beta_{km}(s,t)$ 表示第 $k$ 个基因中位于基因组位置 $t$ 处的 SNP，对第 $m$ 个基因中位于位置 $s$ 处 CpG 位点甲基化水平变异的贡献。

为将 FSEMF 转化为经典结构方程模型，首先汇集全部映射到 $[0,1]$ 区间的基因型函数，并进行 FPCA。随后，用正交归一特征函数 $\psi_l(t)$ 展开 $x_k(t)$，得到
\begin{equation}
x_k(t)=\sum_{l=1}^L\xi_{kl}\psi_l(t).
\tag{5.111}
\end{equation}
将展开式 (5.111) 代入模型 (5.110)，得到
\begin{equation}
\begin{aligned}
&y_1(s)\gamma_{11}+\cdots+y_M(s)\gamma_{M1}
+\sum_{l=1}^L\xi_{1l}b_{l1}^{(1)}(s)\\
&\hspace{12em}+\cdots+\sum_{l=1}^L\xi_{Kl}b_{lK}^{(1)}(s)+e_1(s)=0,\\
&\hspace{14em}\vdots\\
&y_1(s)\gamma_{1M}+\cdots+y_M(s)\gamma_{MM}
+\sum_{l=1}^L\xi_{1l}b_{l1}^{(M)}(s)\\
&\hspace{12em}+\cdots+\sum_{l=1}^L\xi_{Kl}b_{lK}^{(M)}(s)+e_M(s)=0,
\end{aligned}
\tag{5.112}
\end{equation}
其中
\begin{equation}
b_{lk}^{(m)}(s)=\int_T\beta_{km}(s,t)\psi_l(t)\,dt,
\quad m=1,\ldots,M,\quad l=1,\ldots,L,\quad k=1,\ldots,K.
\tag{5.113}
\end{equation}

接着，汇集全部甲基化函数并进行 FPCA，得到一组特征函数 $\phi_j(s)$（$j=1,\ldots,J$）。然后，按特征函数或其他基函数 $\phi_j(s)$ 展开 $y_i(s)$、$b_{lk}^{(m)}(s)$ 和 $e_i(s)$：
\begin{equation}
\begin{aligned}
y_i(s)&=\sum_{j=1}^Jy_{ij}\phi_j(s),\\
b_{lk}^{(m)}(s)&=\sum_{j=1}^J\eta_{lkm}^{(j)}\phi_j(s),\\
e_i(s)&=\sum_{j=1}^J\varepsilon_{ij}\phi_j(s).
\end{aligned}
\tag{5.114}
\end{equation}
将展开式 (5.114) 代入式 (5.112)，得到
\begin{equation}
\begin{aligned}
&y_{1j}\gamma_{11}+\cdots+y_{Mj}\gamma_{M1}
+\sum_{l=1}^L\xi_{1l}\eta_{l11}^{(j)}
+\cdots+\sum_{l=1}^L\xi_{Kl}\eta_{lK1}^{(j)}+\varepsilon_{1j}=0,\\
&\hspace{14em}\vdots\\
&y_{1j}\gamma_{1M}+\cdots+y_{Mj}\gamma_{MM}
+\sum_{l=1}^L\xi_{1l}\eta_{l1M}^{(j)}
+\cdots+\sum_{l=1}^L\xi_{Kl}\eta_{lKM}^{(j)}+\varepsilon_{Mj}=0,\\
&\hfill j=1,\ldots,J.
\end{aligned}
\tag{5.115}
\end{equation}

对于第 $i$ 名个体，按特征函数 $\phi_j(s)$ 和 $\psi_l(t)$ 展开 $y_{ni}(s)$ 与 $x_{nk}(t)$（$n=1,\ldots,T$）：
\[
y_{ni}(s)=\sum_{j=1}^Ly_{nij}\phi_j(s),\qquad
x_{nk}(t)=\sum_{l=1}^L\xi_{nkl}\psi_l(t),\quad n=1,\ldots,T.
\]
定义甲基化函数的第 $j$ 个分量展开系数矩阵 $Y^{(j)}$、基因型函数的展开系数矩阵 $X$、连接甲基化基因的结构参数矩阵 $\Gamma$、连接基因与甲基化基因的结构参数矩阵 $B^{(j)}$，以及误差矩阵 $E^{(j)}$，如下：
\[
Y^{(j)}=\begin{bmatrix}y_{11j}&y_{12j}&\cdots&y_{1Mj}\\\vdots&\vdots&&\vdots\\y_{T1j}&y_{T2j}&\cdots&y_{TMj}\end{bmatrix}
=[y_1^{(j)},y_2^{(j)},\ldots,y_M^{(j)}],
\]
\[
X=\begin{bmatrix}
\xi_{111}&\cdots&\xi_{11L}&\cdots&\xi_{1K1}&\cdots&\xi_{1KL}\\
\vdots&\ddots&\vdots&\reflectbox{$\ddots$}&\vdots&\ddots&\vdots\\
\xi_{T11}&\cdots&\xi_{T1L}&\cdots&\xi_{TK1}&\cdots&\xi_{TKL}
\end{bmatrix}=[\xi_1,\ldots,\xi_K],
\]
\[
\xi_k=\begin{bmatrix}\xi_{1k1}&\cdots&\xi_{1kL}\\\vdots&\vdots&\vdots\\\xi_{TK1}&\cdots&\xi_{TkL}\end{bmatrix},
\]
\[
E^{(j)}-\begin{bmatrix}\varepsilon_{11j}&\varepsilon_{12j}&\cdots&\varepsilon_{1Mj}\\\vdots&\vdots&&\vdots\\\varepsilon_{T1j}&\varepsilon_{T2j}&\cdots&\varepsilon_{TMj}\end{bmatrix}
=[\varepsilon_1^{(j)},\varepsilon_2^{(j)},\ldots,\varepsilon_M^{(j)}],
\]
\[
\Gamma=\begin{bmatrix}r_{11}&r_{12}&\cdots&r_{1M}\\\vdots&\vdots&&\vdots\\r_{M1}&r_{M2}&\cdots&r_{MM}\end{bmatrix}
=[\Gamma_1,\Gamma_2,\ldots,\Gamma_M],
\]
\[
B^{(j)}=\begin{bmatrix}
\eta_{111}^{(j)}&\cdots&\eta_{11M}^{(j)}\\
\vdots&\cdots&\vdots\\
\eta_{L11}^{(j)}&\cdots&\eta_{L1M}^{(j)}\\
\vdots&\cdots&\vdots\\
\eta_{1K1}^{(j)}&\cdots&\eta_{1KM}^{(j)}\\
\vdots&\cdots&\vdots\\
\eta_{LK1}^{(j)}&\cdots&\eta_{LKM}^{(j)}
\end{bmatrix}=[\eta_1^{(j)},\ldots,\eta_M^{(j)}].
\]
式 (5.115) 中的结构方程可写为矩阵形式：
\begin{equation}
Y^{(j)}\Gamma+XB^{(j)}+E^{(j)}=0,\qquad j=1,\ldots,J.
\tag{5.116}
\end{equation}
式 (5.116) 中由第 $j$ 个分量构成的第 $i$ 个方程为
\begin{equation}
y_i^{(j)}=Y_{-i}^{(j)}\gamma_i+\sum_{l=1}^L\xi_{1l}\eta_{l1i}^{(j)}
+\cdots+\sum_{l=1}^L\xi_{Kl}\eta_{lKi}^{(j)}+e_i^{(j)}.
\tag{5.117}
\end{equation}
令
\[
\eta_i^{(j)}=\begin{bmatrix}\eta_{1i}^{(j)}\\\vdots\\\eta_{Ki}^{(j)}\end{bmatrix},\qquad
\eta_{ki}^{(j)}=\begin{bmatrix}\eta_{1ki}^{(j)}\\\vdots\\\eta_{Lki}^{(ij}\end{bmatrix},\qquad
\Delta_i=\begin{bmatrix}\gamma_i\\\eta_i^{(1)}\\\vdots\\\eta_i^{(J)}\end{bmatrix}.
\]
于是，式 (5.117) 可化为
\begin{equation}
y_i^{(j)}=Y_{-i}^{(j)}\gamma_i+X\eta_i^{(j)}+e_i^{(j)},\qquad j=1,\ldots,J,
\tag{5.118}
\end{equation}
其中，$\gamma_i=[\gamma_{1i},\ldots,\gamma_{Mi}]^T$。

在式 (5.118) 两边同乘 $X^T$，得到
\begin{equation}
X^Ty_i^{(j)}=X^TY_{-i}^{(j)}\gamma_i+X^TX\eta_i^{(j)}+X^Te_i^{(j)},\quad j=1,\ldots,J.
\tag{5.119}
\end{equation}

\subsubsection{用测序数据估计基因型—甲基化网络的稀疏 FSEMF}

与第 5.5.2.3 节类似，为建立稀疏 FSEMF，首先定义用于拟合式 (5.119) 的目标函数：
\begin{equation}
\begin{aligned}
f(\Delta_i)=\frac12\sum_{j=1}^J
&\left(X^Ty_i^{(j)}-X^TY_{-i}^{(j)}\gamma_i-X^TX\eta_i^{(j)}\right)^T
\left[\sigma_{ii}^{(j)}X^TX\right]^{-1}\\
&\left(X^Ty_i^{(j)}-X^TY_{-i}^{(j)}\gamma_i-X^TX\eta_i^{(j)}\right).
\end{aligned}
\tag{5.120}
\end{equation}
我们用 $l_1$ 惩罚连接甲基化基因的参数 $\gamma_i$，用组套索惩罚连接基因与甲基化基因的参数 $\eta_i^{(j)}$。为此，建立稀疏 FSEMF：
\begin{equation}
\min\quad f(\Delta_i)+\lambda\left[\lVert Z_{\gamma i}\rVert_1
+\sqrt{JL}\sum_{k=1}^K\lVert Z_{ki}\rVert_2\right]
\tag{5.121}
\end{equation}
约束为 $\Delta_i-Z_i=0$，其中
\[
\widetilde\Delta_i=\begin{bmatrix}\widetilde\Delta_{\gamma i}\\\widetilde\Delta_{1i}\\\vdots\\\widetilde\Delta_{Ki}\end{bmatrix},\quad
\widetilde\Delta_{\gamma i}=\gamma_i,\quad
\widetilde\Delta_{ki}=\begin{bmatrix}\widetilde\Delta_{1ki}\\\vdots\\\widetilde\Delta_{Lki}\end{bmatrix},
\]
\[
\widetilde\Delta_{lki}=\begin{bmatrix}\widetilde\Delta_{lki}^{(1)}\\\vdots\\\widetilde\Delta_{lki}^{(J)}\end{bmatrix},\qquad
\widetilde\Delta_{lki}^{(j)}=\eta_{lki}^{(j)},\qquad Z_{lki}^{(j)}=\eta_{lki}^{(j)},
\]
\[
Z_i^{(j)}=\begin{bmatrix}Z_{1i}^{(j)}\\\vdots\\Z_{Ki}^{(j)}\end{bmatrix},\quad
Z_{ki}=\begin{bmatrix}Z_{1ki}\\\vdots\\Z_{Lki}\end{bmatrix},\quad
Z_{lki}=\begin{bmatrix}Z_{lki}^{(1)}\\\vdots\\Z_{lki}^{(J)}\end{bmatrix},\quad
Z_i=\begin{bmatrix}Z_{\gamma i}\\Z_{1i}\\\vdots\\Z_{Ki}\end{bmatrix},
\]
且 $Z_{\gamma i}=\gamma_i$。

为求解优化问题 (5.121)，构造增广拉格朗日函数
\begin{equation}
\begin{aligned}
L_\rho(\Delta_i,Z_i,\mu)={}&f(\Delta_i)+\lambda\lVert Z_{\gamma i}\rVert_1
+\lambda\sum_{k=1}^K\sqrt J\lVert Z_{ki}\rVert_2\\
&+\frac\rho2\lVert\Delta_i-Z_i+\mu\rVert_2^2-\frac\rho2\lVert\mu\rVert_2^2.
\end{aligned}
\tag{5.122}
\end{equation}
求解该优化问题的交替方向乘子法（ADMM）包含以下迭代：
\begin{equation}
\Delta_i^{(l+1)}:=\argmin_{\Delta_i}L_\rho(\Delta_i,Z_i^{(l)},\mu^{(l)}),
\tag{5.123}
\end{equation}
\begin{equation}
Z_i^{(l+1)}:=\argmin_{Z_i}L_\rho(\Delta_i^{(l+1)},Z_i,\mu^{(l)}),
\tag{5.124}
\end{equation}
\begin{equation}
\mu^{(l+1)}:=\mu^{(l)}+\Delta_i^{(l+1)}-Z_i^{(l+1)}.
\tag{5.125}
\end{equation}
或者写为
\begin{equation}
\Delta_i^{(l+1)}:=\argmin_{\Delta_i}\left(f(\Delta_i)
+\frac\rho2\lVert\Delta_i-Z_i^{(l)}+u^{(l)}\rVert_2^2\right),
\tag{5.126}
\end{equation}
\begin{equation}
\begin{aligned}
Z_i^{(l+1)}:=\argmin_{Z_i}\Biggl(&\lambda\lVert Z_{\gamma i}\rVert_1+\lambda
+\sum_{k=1}^K\sqrt J\lVert Z_{ki}\rVert_2\\
&+\frac\rho2\lVert\Delta_i^{(l+1)}-Z_i+\mu^{(l)}\rVert_2^2\Biggr),
\end{aligned}
\tag{5.127}
\end{equation}
\begin{equation}
\mu^{(l+1)}:=\mu^{(l)}+\Delta_i^{(l+1)}-Z_i^{(l+1)}.
\tag{5.128}
\end{equation}

现在求解极小化问题 (5.126)。该问题仅涉及可微函数，直接运用微积分即可求解。令
\[
\mu=\begin{bmatrix}\mu_\gamma\\\mu_1\\\vdots\\\mu_J\end{bmatrix}
\qquad\text{或}\qquad
\mu=\begin{bmatrix}\mu_\gamma\\\mu_{1i}\\\vdots\\\mu_{Ki}\end{bmatrix}.
\]
令 $L_\rho$ 的偏导数为零，得到
\begin{equation}
\begin{aligned}
\frac{\partial L_\rho}{\partial r_i}
={}&-\sum_{j=1}^J(Y_{-i}^{(j)})^TX(\sigma_{ii}^{(j)}X^TX)^{-1}\\
&\quad\cdot(X^Ty_i^{(j)}-X^TY_{-i}^{(j)}\gamma_i-X^TX\eta_i^{(j)})
+\rho(\gamma_i-Z_{\gamma i}^{(l)}+\mu_\gamma^{(l)})=0,\\
\frac{\partial L_\rho}{\partial\eta_i^{(j)}}
={}&-X^TX(\sigma_{ii}^{(j)}X^TX)^{-1}
(X^Ty_i^{(j)}-X^TY_{-i}^{(j)}\gamma_i-X^TX\eta_i^{(j)})\\
&+\rho(\eta_i^{(j)}-Z_i^{(j)(l)}+\mu_j^{(l)})=0,\quad j=1,\ldots,J.
\end{aligned}
\tag{5.129}
\end{equation}
式 (5.129) 可转化为
\begin{equation}
\begin{aligned}
(A_1+\rho I)\gamma_i+\sum_{j=1}^JB_1(j)\eta_i^{(j)}&=D_1,\\
A_2(j)\gamma_i+[B_2(j)+\rho I]\eta_i^{(j)}&=D_2(j),\quad j=1,\ldots,J,
\end{aligned}
\tag{5.130}
\end{equation}
其中
\[
\begin{aligned}
A_1&=\sum_{j=1}^J(Y_{-i}^{(j)})^TX(\sigma_{ii}^{(j)}X^TX)^{-1}X^TY_{-i}^{(j)},\\
A_2(j)&=X^TX(\sigma_{ii}^{(j)}X^TX)^{-1}X^TY_{-i}^{(j)},\\
B_1(j)&=(Y_{-i}^{(j)})^TX(\sigma_{ii}^{(j)}X^TX)^{-1}X^TX,\\
B_2(j)&=X^TX(\sigma_{ii}^{(j)}X^TX)^{-1}X^TX,\\
D_1&=\sum_{j=1}^J(Y_{-i}^{(j)})^TX(\sigma_{ii}^{(j)}X^TX)^{-1}X^Ty_i^{(j)}
+\rho(Z_{i\gamma}^{(l)}-\mu_\gamma^{(l)}),\\
D_2(j)&=X^TX(\sigma_{ii}^{(j)}X^TX)^{-1}X^Ty_i^{(j)}
+\rho(Z_i^{(j)(l)}-\mu_j^{(l)}).
\end{aligned}
\]
式 (5.130) 的解为
\begin{equation}
\begin{aligned}
\widehat\gamma_i&=\left[A_1-\sum_{j=1}^JB_1(j)(B_2(j)+\rho I)^{-1}A_2(j)\right]^{-1}D,\\
D&=D_1-\sum_{j=1}^JB_1(j)(B_2(j)+\rho I)^{-1}D_2(j),\\
\eta_i^{(j)}&=[B_2(j)+\rho I]^{-1}[D_2(j)-A_2(j)\widehat\gamma_i],\quad j=1,\ldots,J.
\end{aligned}
\tag{5.131}
\end{equation}

优化问题 (5.127) 不可微。尽管式中前两项不可微，仍可利用次微分运算，得到简单的闭式解。先考虑 $\lVert Z_{\gamma i}\rVert_1$ 的广义梯度。令 $\Phi_i^{(m)}$ 为向量 $Z_{\gamma i}$ 第 $m$ 个分量的广义导数：
\[
\Phi_i^{(m)}=\begin{cases}
1,&Z_{\gamma i}^{(m)}=\gamma_i^{(m)}>0,\\
[-1,1],&Z_{\gamma i}^{(m)}=\gamma_i^{(m)}=0,\\
-1,&Z_{\gamma i}^{(m)}=\gamma_i^{(m)}<0.
\end{cases}
\]
令 $\Phi_i=[\Phi_i^{(1)},\ldots,\Phi_i^{(M)}]^T$。于是
\[
\frac\lambda\rho\Phi_i+Z_{\gamma i}=\Delta_{\gamma i}^{(l+1)}+\mu_\gamma^{(l)},
\]
从而有
\begin{equation}
Z_{\lambda i}^{(l+1)}=\sgn(\Delta_{\gamma i}^{(l+1)}+\mu_\gamma^{(l)})
\left(|\Delta_{\gamma i}^{(l+1)}+\mu_\gamma^{(l)}|-\frac\lambda\rho\right)_+,
\tag{5.132}
\end{equation}
其中
\[
|x|_+=\begin{cases}x,&x\geq0,\\0,&x<0.\end{cases}
\]
接着考虑组套索。广义梯度 $\partial L_\rho/\partial Z_{ki}$ 为
\begin{equation}
\frac{\partial L_\rho}{\partial Z_{ki}}
=\lambda\sqrt J\,s+\rho\left(Z_{ki}-(\Delta_{ki}^{(l+1)}-\mu_{ki}^{(l)})\right)=0,
\tag{5.133}
\end{equation}
其中
\[
s=\begin{cases}
0,&\lVert\Delta_{ki}^{(l+1)}-\mu_{ki}^{(l)}\rVert_2<\dfrac{\lambda\sqrt J}\rho,\\[4pt]
\dfrac{Z}{\lVert Z_{ki}\rVert_2},&\lVert\Delta_{ki}^{(l+1)}-\mu_{ki}^{(l)}\rVert_2\geq\dfrac{\lambda\sqrt J}\rho.
\end{cases}
\]
式 (5.133) 可化为
\begin{equation}
\left(1+\frac{\lambda\sqrt J}{\rho\lVert Z_{ki}\rVert_2}\right)Z_{ki}
=\Delta_{ki}^{(l+1)}-\mu_{ki}^{(l)},
\tag{5.134}
\end{equation}
或者
\begin{equation}
\left(1+\frac{\lambda\sqrt J}{\rho\lVert Z_{ki}\rVert_2}\right)\lVert Z_{ki}\rVert_2
=\lVert\Delta_{ki}^{(l+1)}-\mu_{ki}^{(l)}\rVert_2.
\tag{5.135}
\end{equation}
解式 (5.135)，得到
\begin{equation}
\lVert Z_{ki}\rVert_2=\lVert\Delta_{ki}^{(l+1)}-\mu_{ki}^{(l)}\rVert_2
-\frac{\lambda\sqrt J}\rho.
\tag{5.136}
\end{equation}
将式 (5.136) 代入式 (5.135)，得到
\begin{equation}
\left(1+\frac{\lambda\sqrt J}{\rho\lVert Z_{ki}\rVert_2}\right)
=\frac{1}{1-\dfrac{\lambda\sqrt J}{\rho\lVert\Delta_{ki}^{(l+1)}-\mu_{ki}^{(l)}\rVert_2}}.
\tag{5.137}
\end{equation}
结合式 (5.134)、(5.135) 和 (5.137)，得到
\begin{equation}
Z_{ki}^{(l+1)}=\left(-\frac{\lambda\sqrt J}{\rho\lVert\Delta_{ki}^{(l+1)}-\mu_{ki}^{(l)}\rVert_2}\right)_+
(\Delta_{ki}^{(l+1)}-\mu_{ki}^{(l)}),\quad k=1,\ldots,K.
\tag{5.138}
\end{equation}
最后，汇总上述各式，得到结论 5.10。

\par\medskip\Needspace{5\baselineskip}
\noindent\textbf{结论 5.10：将甲基化水平与基因型函数均作为函数值性状的基因型—甲基化因果网络构建算法}\par
对于 $i=1,\ldots,M$：
\begin{enumerate}[label=步骤 \arabic*：,leftmargin=4.8em]
\item 初始化。
\[
\begin{aligned}
B_1(j)&=(Y_{-i}^{(j)})^TX(\sigma_{ii}^{(j)}X^TX)^{-1}X^TX,\\
B_2(j)&=X^TX(\sigma_{ii}^{(j)}X^TX)^{-1}X^TX,\\
\widetilde D_1&=\sum_{j=1}^J(Y_{-i}^{(j)})^TX(\sigma_{ii}^{(j)}X^TX)^{-1}X^Ty_i^{(j)},\\
\widetilde D_2(j)&=X^TX(\sigma_{ii}^{(j)}X^TX)^{-1}X^Ty_i^{(j)}.
\end{aligned}
\]
\[
\mu^0:=0,
\]
\[
\begin{aligned}
\gamma_i^{(0)}&=\left[A_1-\sum_{j=1}^JB_1(j)(B_2(j)+\rho I)^{-1}A_2(j)\right]^{-1}\widetilde D,\\
\widetilde D&=\widetilde D_1-\sum_{j=1}^JB_1(j)(B_2(j)+\rho I)^{-1}\widetilde D_2(j),\\
\eta_i^{(0)(j)}&=[B_2(j)+\rho I]^{-1}[D_2(j)-A_2(j)\widehat\gamma_i^{(0)}],\quad j=1,\ldots,J,
\end{aligned}
\]
\[
\Delta_i^{(0)}=\begin{bmatrix}\gamma_i^{(0)}\\\eta_i^{(0)(1)}\\\vdots\\\eta_i^{(0)(J)}\end{bmatrix},\qquad
Z_i^{(0)}:==\Delta_i^{(0)},
\]
\[
\widetilde\Delta_i^{(0)}=\begin{bmatrix}\widetilde\Delta_{\gamma i}^{(0)}\\\widetilde\Delta_{1i}^{(0)}\\\vdots\\\widetilde\Delta_{Ki}^{(0)}\end{bmatrix},\quad
\widetilde\Delta_{\gamma i}^{(0)}=\widehat\gamma_i^{(0)},\quad
\widetilde\Delta_{ki}^{(0)}=\begin{bmatrix}\widetilde\Delta_{ki}^{(0)(1)}\\\vdots\\\widetilde\Delta_{ki}^{(0)(J)}\end{bmatrix},
\]
\[
\widetilde\Delta_{ki}^{(0)(j)}=\eta_{ki}^{(0)(j)},\qquad
\mu^{(0)}=\begin{bmatrix}\mu_\gamma^{(0)}\\\mu_1^{(0)}\\\vdots\\\mu_J^{(0)}\end{bmatrix}.
\]
执行步骤 2、3 和 4，直至收敛。
\item
\[
\begin{aligned}
A_1&=\sum_{j=1}^J(Y_{-i}^{(j)})^TX(\sigma_{ii}^{(j)}X^TX)^{-1}X^TY_{-i}^{(j)},\\
A_2(j)&=X^TX(\sigma_{ii}^{(j)}X^TX)^{-1}X^TY_{-i}^{(j)},\\
B_1(j)&=(Y_{-i}^{(j)})^TX(\sigma_{ii}^{(j)}X^TX)^{-1}X^TX,\\
B_2(j)&=X^TX(\sigma_{ii}^{(j)}X^TX)^{-1}X^TX,\\
D_1&=\sum_{j=1}^J(Y_{-1}^{(j)})^TX(\sigma_{ii}^{(j)}X^TX)^{-1}X^Ty_i^{(j)}
+\rho(Z_{\gamma i}^{(l)}-\mu_\gamma^{(l)}),\\
D_2(j)&=X^TX(\sigma_{ii}^{(j)}X^TX)^{-1}X^Ty_i^{(j)}
+\rho(Z_i^{(j)(l)}-\mu_j^{(l)}).
\end{aligned}
\]
\[
\begin{aligned}
\widehat\gamma_i^{(l+1)}&=\left[A_1-\sum_{j=1}^JB_1(j)(B_2(j)+\rho I)^{-1}A_2(j)\right]^{-1}D,\\
D&=D_1-\sum_{j=1}^JB_1(j)(B_2(j)+\rho I)^{-1}D_2(j),\\
\eta_i^{(j)(l+1)}&=[B_2(j)+\rho I]^{-1}[D_2(j)-A_2(j)\widehat\gamma_i^{(l+1)}],\quad j=1,\ldots,J,
\end{aligned}
\]
\[
\Delta_i^{(l+1)}=\begin{bmatrix}\widehat\gamma_i^{(l+1)}\\\eta_i^{(1)(l+1)}\\\vdots\\\eta_i^{(J)(l+1)}\end{bmatrix},
\]
\[
\widetilde\Delta_i^{(l+1)}=\begin{bmatrix}\widetilde\Delta_{\gamma i}^{(l+1)}\\\widetilde\Delta_{1i}^{(l+1)}\\\vdots\\\widetilde\Delta_{Ki}^{(l+1)}\end{bmatrix},\quad
\widetilde\Delta_{\gamma i}^{(l+1)}=\widehat\gamma_i^{(l+1)},\quad
\widetilde\Delta_{ki}^{(l+1)}=\begin{bmatrix}\widetilde\Delta_{ki}^{(l+1)(1)}\\\vdots\\\widetilde\Delta_{ki}^{(l+1)(J)}\end{bmatrix},
\]
\[
\widetilde\Delta_{ki}^{(l+1)(j)}=\eta_{ki}^{(j)(l+1)}.
\]
\item
\[
\widetilde Z_{\gamma i}^{(l+1)}=\sgn(\widetilde\Delta_{\gamma i}^{(l+1)}+\mu_\gamma^{(l)})
\left(|\widetilde\Delta_{\gamma i}^{(l+1)}+\mu_\gamma^{(l)}|-\frac\lambda\rho\right)_+,
\]
其中
\[
|x|_+=\begin{cases}x,&x\geq0,\\0,&x<0.\end{cases}
\]
\[
\widetilde Z_{ki}^{(l+1)}=\left(-\frac{\lambda\sqrt J}{\rho\lVert\widetilde\Delta_{ki}^{(l+1)}-\mu_{ki}^{(l)}\rVert_2}\right)_+
(\widetilde\Delta_{ki}^{(l+1)}-\mu_{ki}^{(l)}),\quad k=1,\ldots,K.
\]
\[
\widetilde Z_i^{(l+1)}=\begin{bmatrix}\widetilde Z_{\gamma i}^{(l+1)}\\\widetilde Z_{1i}^{(l+1)}\\\vdots\\\widetilde Z_{Ki}^{(l+1)}\end{bmatrix},\qquad
\widetilde Z_{ki}^{(l+1)}=\begin{bmatrix}\widetilde Z_{ki}^{(1)(l+1)}\\\vdots\\\widetilde Z_{ki}^{(J)(l+1)}\end{bmatrix}.
\]
更新 $Z^{(l+1)}$：
\[
Z_{ki}^{(j)(l+1)}=\widetilde Z_{ki}^{(j)(l+1)},\qquad
Z_{\gamma i}^{(l+1)}=\widetilde Z_{\gamma i}^{(l+1)},
\]
\[
Z_i^{(j)(l+1)}=\begin{bmatrix}Z_{1i}^{(j)(l+1)}\\\vdots\\Z_{Ki}^{(j)(l+1)}\end{bmatrix},\qquad
Z_i^{(l+1)}=\begin{bmatrix}Z_{\gamma i}^{(l+1)}\\Z_i^{(1)(l+1)}\\\vdots\\Z_i^{(J)(l+1)}\end{bmatrix}.
\]
\item
\[
\widetilde\mu^{(l+1)}=\widetilde\mu_i^{(l)}+(\widetilde\Delta_i^{(l+1)}-\widetilde Z_i^{(l+1)}),
\]
其中
\[
\widetilde\mu^{(l+1)}=\begin{bmatrix}\widetilde\mu_\gamma^{(l+1)}\\\widetilde\mu_1^{(l+1)}\\\vdots\\\widetilde\mu_K^{(l+1)}\end{bmatrix},\qquad
\widetilde\mu_k^{(l+1)}=\begin{bmatrix}\widetilde\mu_k^{(1)(l+1)}\\\vdots\\\widetilde\mu_k^{(J)(l+1)}\end{bmatrix},\quad k=1,\ldots,K.
\]
更新 $\mu^{(l+1)}$：
\[
\mu^{(l+1)}=\begin{bmatrix}\mu_\gamma^{(l+1)}\\\mu_1^{(l+1)}\\\vdots\\\mu_J^{(l+1)}\end{bmatrix},\qquad
\mu_\gamma^{(l+1)}=\widetilde\mu_\gamma^{(l+1)},
\]
\[
\mu_j^{(l+1)}=\begin{bmatrix}\mu_1^{(j)(l+1)}\\\vdots\\\mu_K^{(j)(l+1)}\end{bmatrix},\qquad
\mu_k^{(j)(l+1)}=\widetilde\mu_k^{(j)(l+1)}.
\]
返回步骤 2。
\end{enumerate}

\section*{软件包}
\phantomsection\addcontentsline{toc}{section}{软件包}
\markright{软件包}

软件包 minfi 用于分析 Human Methylation450（“450k”）和 EPIC 平台的 Illumina DNA 甲基化芯片数据，可从 \url{https://github.com/kasperdanielhansen/minfi} 下载。ChAMP 是一个集成分析流程，发表于 2014 年，并于 2017 年修订；下载地址为 \url{https://bioconductor.org/packages/release/bioc/vignettes/ChAMP/inst/doc/ChAMP.html}。

Illumina 提供支持数据整合的数据分析软件，例如，将基因分型与基因表达数据结合进行 eQTL 分析，或与甲基化数据结合进行 mQTL 分析（技术说明：Illumina® 系统与软件——Illumina 数据的 QTL 分析软件工具）。“meQTL 定位分析实用指南”（meQTL mapping analysis cookbook）也提供了 meQTL 分析软件。
